\chapter*{Preface}
\markboth{PREFACE}{PREFACE}

The first German edition of this volume was published by Julius Springer, Berlin, in 1924. A second edition, revised and improved with the help of K. O. Friedrichs, R. Luneburg, F. Rellich, and other unselfish friends, followed in 1930. The second volume appeared in 1938. In the meantime I had been forced to leave Germany and was fortunate and grateful to be given the opportunities open in the United States. During the Second World War the German book became unavailable and later was even suppressed by the National Socialist rulers of Germany. Thus the survival of the book was secured when the United States Government seized the copyright and licensed a reprint issued by Interscience Publishers, New York. Such a license also had to be obtained from the Alien Property Custodian for the present English edition.

This edition follows the German original fairly closely but contains a large number of additions and modifications. I have had to postpone a plan to completely rewrite and modernize the book in collaboration with K. O. Friedrichs, because the pressure for publication of an English ``Courant-Hilbert'' has become irresistible. Even so, it is hoped that the work in its present form will be useful to mathematicians and physicists alike, as the numerous demands from all sides seem to indicate.

The objective of the book can still today be expressed almost as in the preface to the first German edition. ``Since the seventeenth century, physical intuition has served as a vital source for mathematical problems and methods. Recent trends and fashions have, however, weakened the connection between mathematics and physics; mathematicians, turning away from the roots of mathematics in intuition, have concentrated on refinement and emphasized the postulational side of mathematics, and at times have overlooked the unity of their science with physics and other fields. In many cases, physicists have ceased to appreciate the attitudes of mathematicians. This rift is unquestionably a serious threat to science as a whole; the broad stream of scientific development may split into smaller and smaller rivulets and dry out. It seems therefore important to direct our efforts toward reuniting divergent trends by clarifying the common features and interconnections of many distinct and diverse scientific facts. Only thus can the student attain some mastery of the material and the basis be prepared for further organic development of research.

``The present work is designed to serve this purpose for the field of mathematical physics. Mathematical methods originating in problems of physics are developed and the attempt is made to shape results into unified mathematical theories. Completeness is not attempted, but it is hoped that access to a rich and important field will be facilitated by the book.

``The responsibility for the present book rests with me. Yet the name of my teacher, colleague, and friend, D. Hilbert, on the title page seems justified by the fact that much material from Hilbert's papers and lectures has been used, as well as by the hope that the book expresses some of Hilbert's spirit, which has had such a decisive influence on mathematical research and education.''

I am greatly indebted to many helpers in all phases of the task of preparing this edition: to Peter Ceike, Ernest Courant, and Anneli Lax, who provided most of the first draft of the translation; to Hanan Rubin and Herbert Kranzer, who have given constructive criticism; to Wilhelm Magnus, who is responsible for the appendix to Chapter VII; and to Natascha Artin and Lucile Gardner, who carried the burden of the editorial work. Most cordial thanks also are due to Interscience Publishers for their patient and helpful attitude and to my old friend and publisher, Dr. Ferdinand Springer in Heidelberg, the great pioneer of modern scientific publishing, for his sympathetic understanding of the situation, which has so greatly changed since the old days of our close cooperation.

\vspace{1.8em}
\begin{flushright}
R. COURANT\\
NEW ROCHELLE, NEW YORK\\
JUNE 1953
\end{flushright}
\clearpage
